\subsection{初始分歧指数}

定义 4. 设 G 是一个全序交换群，H 是 G 的有限指数子群。G 的由严格正元素组成且包含 H 的所有元素 \textgreater0 的优势子集的数目，称为 H 在 G 中的初始指数，记作 \(\varepsilon(\mathbf{G}, \mathbf{H})\)。

根据下述命题，这个初始指数是自然数：

命题 3. 在定义 4 的假设下，若 G 的严格正元素集没有最小元素，则对所有 H 都有 \(\varepsilon(G, H) = 1\)。如果 G 存在最小的 \(>0\) 元素，并记 G' 为由它生成的子群，则

\[
\varepsilon (\mathrm{G}, \mathrm{H}) = (\mathrm{G} ^ {\prime}: (\mathrm{G} ^ {\prime} \cap \mathrm{H})).
\]

在第一种情形中，令 x 为 G 中的一个元素 \textgreater0。满足 \(y \in G\) 且 0 \textless{} y \textless{} x 的集合是无限的，因而其中存在两个不同但模 H 同余的元素；它们的差是 H 中满足 0 \textless{} z \textless{} x 的元素。因此，每个包含 H 的全部严格正元素的优势子集都包含 x，进而包含 G 的所有元素 \textgreater0。

在第二种情形中，令 \(x\) 为最小的 \(>0\) 元素（在 \(G\) 中），并令 \(n\) 为最小的 \(>0\) 整数，使 \(nx \in H\)。显然 \(n = (G': (G' \cap H))\)。另一方面，记 \(M(y)\) 为满足 \(z \in G\) 且 \(y \leqslant z\)（\(y \in G\)）的 z 的集合，立即可见，定义 4 的优势集恰为 \(M(x), M(2x), \ldots, M(nx)\)。

推论。初始指数 \(\varepsilon(\mathbf{G},\mathbf{H})\) 整除指数 (G: H)，当 G 同构于 Z 时它等于该指数。

特别地，\(\varepsilon(G, H) \leqslant (G: H)\)。

定义 5. 设 K 为域，L 为 K 的有限扩张，w 为 L 上的赋值，v 为其在 K 上的限制，\(\Gamma_{w}\) 和 \(\Gamma_{v}\) 为它们的阶群。\(\Gamma_{v}\) 在 \(\Gamma_{w}\) 中的初始指数称为 w 关于 v（或关于 K）的初始分歧指数，记作 \(\varepsilon(w/v)\)。

由上述推论，\(\varepsilon(w/v)\) 整除 \(e(w/v)\)，在离散赋值的情形下二者相等。

命题 4. 在定义 5 的假设下，设 A 和 m（分别地，A' 和 m'）为赋值 v（分别地，w）的环和理想。则

\[
[ \mathrm{A} ^ {\prime} / \mathrm{mA} ^ {\prime}: \mathrm{A} / \mathrm{m} ] = \varepsilon (w / v) f (w / v).
\]

包含 mA' 且不同于 A' 的 A' 的理想，与由 \(\Gamma_{w}\) 的元素 \textgreater0 组成并包含 \textgreater0 的 \(\Gamma_{v}\) 的元素的优势子集相对应（\(\S3\)，第 5 节，命题 7 的推论）。因此它们的数目等于 \(\varepsilon(w/v)\)；由于它们关于包含关系构成全序集，这个数目等于商环 A'/mA' 的长度。现在，A' 上长度为 1 的模是 A'/m' 上的一维向量空间，因而是 A 上长度为 \(f(w/v)\) 的模；所以，既然 A'/mA' 作为 A'-模的长度为 \(\varepsilon(w/v)\)，它作为 A-模，也就是作为 A/m-模的长度为 \(\varepsilon(w/v)f(w/v)\)。

